\NormalFont\ShortTitle{Apocalipsa}

{\MT Apocalipsa lui Ioan


\par }\ChapOne{1}{\PP \VerseOne{1}Aceasta este Revelația lui Isus Hristos, pe care i-a dat-o Dumnezeu ca să arate robilor Săi lucrurile care trebuie să se întâmple curând, pe care a trimis-o și a făcut-o cunoscută prin îngerul Său robului Său, Ioan,

\VS{2}care a mărturisit cuvântul lui Dumnezeu și mărturia lui Isus Hristos, despre tot ce a văzut.


\par }{\PP \VS{3}Ferice de cel ce citește și de cei ce ascultă cuvintele proorociei și păzesc cele scrise în ea, căci timpul este aproape.


\par }{\PP \VS{4}Ioan, către cele șapte adunări care sunt în Asia: Har vouă și pace de la Dumnezeu, care este, care era și care va veni, și de la cele șapte duhuri care sunt înaintea tronului Lui,

\VS{5}și de la Isus Hristos, martorul credincios, întâiul născut dintre cei morți și conducătorul regilor pământului. Celui care ne iubește și ne-a spălat de păcatele noastre cu sângele său —

\VS{6}și ne-a făcut să fim un regat, preoți ai Dumnezeului și Tatălui său — lui să fie gloria și stăpânirea în vecii vecilor. Amin.


\par }{\PP \VS{7}Iată, El vine cu norii și orice ochi Îl va vedea, chiar și cei ce L-au străpuns. Toate semințiile pământului îl vor plânge. Chiar și așa, Amin.


\par }{\PP \VS{8}{\WJ{“Eu sunt Alfa și Omega,}} zice Domnul Dumnezeu,
{\WJ{Cel ce este și Cel ce era și Cel ce va veni, Cel Atotputernic.”
}}


\par }{\PP \VS{9}Eu, Ioan, fratele vostru și partenerul vostru în asuprire, în Împărăție și în perseverența în Hristos Isus, am fost pe insula numită Patmos, din cauza Cuvântului lui Dumnezeu și a mărturiei lui Isus Hristos.

\VS{10}Eram în Duhul Sfânt în ziua Domnului și am auzit în spatele meu un glas puternic, ca o trâmbiță

\VS{11}care spunea:
\FTNT{*}{{\FR 1:11 }{\FT{TR adaugă: “Eu sunt Alfa și Omega, Cel dintâi și Cel de pe urmă”.}}}{\WJ{“Ce vezi, scrie într-o carte și trimite la cele șapte adunări:}}
\FTNT{†}{{\FR 1:11 }{\FT{TR adaugă “care sunt în Asia”}}}{\WJ{la Efes, Smirna, Pergam, Tiatira, Sardes, Filadelfia și Laodiceea”.
}}


\par }{\PP \VS{12}M-am întors să văd vocea care vorbea cu mine. După ce m-am întors, am văzut șapte sfeșnice de aur.

\VS{13}Și printre sfeșnice era unul ca un fiu de om,
\FTNT{‡}{{\FR 1:13 }{\FT{Daniel 7:13
}}}îmbrăcat cu o haină care-i ajungea până la picioare și cu o cingătoare de aur în jurul pieptului.

\VS{14}Capul și părul lui erau albe ca lâna albă, ca zăpada. Ochii lui erau ca o flacără de foc.

\VS{15}Picioarele lui erau ca arama lustruită, ca și cum ar fi fost rafinată într-un cuptor. Glasul lui era ca un glas de ape multe.

\VS{16}Avea șapte stele în mâna dreaptă. Din gura lui ieșea o sabie ascuțită cu două tăișuri. Fața lui era ca soarele care strălucește la maxim.

\VS{17}Când l-am văzut, am căzut la picioarele lui ca un mort.

\par }{\PP Și-a pus mâna dreaptă peste mine și mi-a zis:
{\WJ{“Nu te teme. Eu sunt Cel dintâi și Cel de pe urmă,
}}

\VS{18}{\WJ{și Cel viu. Am fost mort și iată că sunt viu în vecii vecilor. Amin. Eu am cheile Morții și ale Hadesului.
}}\FTNT{§}{{\FR 1:18 }{\FT{sau, iad}}}

\VS{19}{\WJ{Scrieți, așadar, lucrurile pe care le-ați văzut, lucrurile care sunt și lucrurile care se vor întâmpla în viitor.
}}

\VS{20}{\WJ{Taina celor șapte stele pe care le-ai văzut în mâna Mea dreaptă și a celor șapte sfeșnice de aur este aceasta: Cele șapte stele sunt îngerii}}\FTNT{*}{{\FR 1:20 }{\FT{sau, mesageri (aici și oriunde sunt menționați îngerii)}}}{\WJ{celor șapte adunări. Cele șapte sfeșnice sunt cele șapte adunări.
}}



\par }\Chap{2}{\PP \VerseOne{1}{\WJ{“Scrieți îngerului adunării din Efes:
}}

\par }{\PP {\WJ{“Cel ce ține cele șapte stele în mâna dreaptă, cel ce umblă printre cele șapte steluțe de aur spune aceste lucruri:
}}


\par }{\PP \VS{2}{\WJ{“Știu faptele voastre, și osteneala și stăruința voastră, și că nu puteți îngădui pe cei răi, și că ați pus la încercare pe cei ce se numesc apostoli și nu sunt apostoli, și i-ați găsit mincinoși.
}}

\VS{3}{\WJ{Voi aveți perseverență și ați îndurat pentru Numele Meu și nu}}\FTNT{*}{{\FR 2:3 }{\FT{TR adaugă “au muncit și”}}}{\WJ{ați obosit.
}}

\VS{4}{\WJ{Dar am împotriva voastră acest lucru: că v-ați părăsit dragostea dintâi.
}}

\VS{5}{\WJ{Amintește-ți, așadar, de unde ai căzut, pocăiește-te și fă cele dintâi fapte, altfel vin repede la tine și îți voi muta suportul de lampă de la locul lui, dacă nu te pocăiești.
}}

\VS{6}{\WJ{Dar aveți acest lucru: că urâți faptele nicolaeților, pe care le urăsc și eu.
}}

\VS{7}{\WJ{Cine are ureche, să asculte ce spune Duhul Sfânt adunărilor. Celui care va birui îi voi da să mănânce din pomul vieții, care este în Paradisul Dumnezeului meu.
}}


\par }{\PP \VS{8}{\WJ{“Scrieți îngerului adunării din Smirna:
}}

\par }{\PP {\WJ{“Cel dintâi și cel de pe urmă, care a fost mort și a înviat, spune aceste lucruri:
}}


\par }{\PP \VS{9}{\WJ{“Cunosc faptele voastre, asuprirea și sărăcia voastră (dar voi sunteți bogați) și hulele celor ce se zic iudei, dar nu sunt iudei, ci sunt o sinagogă a Satanei.
}}

\VS{10}{\WJ{Nu vă temeți de lucrurile pe care urmează să le suferiți. Iată, diavolul este pe cale să arunce pe unii dintre voi în închisoare, ca să fiți puși la încercare, și veți avea parte de opresiune timp de zece zile. Fiți credincioși până la moarte și vă voi da cununa vieții.
}}

\VS{11}{\WJ{Cine are urechi, să asculte ce spune Duhul Sfânt adunărilor. Cel care va birui nu va fi vătămat de moartea a doua.
}}


\par }{\PP \VS{12}{\WJ{“Scrieți îngerului adunării din Pergam:
}}

\par }{\PP {\WJ{“Cel care are sabia ascuțită cu două tăișuri spune aceste lucruri:
}}


\par }{\PP \VS{13}{\WJ{Știu faptele tale și locul unde locuiești, unde este tronul lui Satana. Voi țineți cu tărie la numele meu și nu ați negat credința mea în zilele lui Antipa, martorul meu, credinciosul meu, care a fost ucis în mijlocul vostru, unde locuiește Satana.
}}

\VS{14}{\WJ{Dar am câteva lucruri împotriva voastră, pentru că aveți acolo unii care țin învățătura lui Balaam, care l-a învățat pe Balac să arunce o piatră de poticnire înaintea copiilor lui Israel, să mănânce lucruri sacrificate idolilor și să comită imoralitate sexuală.
}}

\VS{15}Tot așa și
{\WJ{voi, de asemenea, aveți unii care țin învățătura nicolahilor.
}}\FTNT{†}{{\FR 2:15 }{\FT{TR citește “pe care îl urăsc” în loc de “de asemenea”}}}

\VS{16}{\WJ{Pocăiți-vă, deci, căci altfel vin repede la voi și voi face război împotriva lor cu sabia gurii mele.
}}

\VS{17}{\WJ{Cine are urechi, să asculte ce spune Duhul Sfânt în adunări. Celui care va învinge, îi voi da din mana ascunsă și}}\FTNT{‡}{{\FR 2:17 }{\FT{Manna este o hrană supranaturală, numită astfel după cuvântul ebraic pentru “Ce este?”. A se vedea Exodul 11:7-9.}}}{\WJ{îi voi da o piatră albă, iar pe piatră va fi scris un nume nou, pe care nimeni nu-l cunoaște decât cel care îl primește.
}}


\par }{\PP \VS{18}{\WJ{“Scrieți îngerului adunării din Tiatira:
}}

\par }{\PP {\WJ{“Fiul lui Dumnezeu, care are ochii ca o flacără de foc și picioarele ca arama lustruită, spune următoarele:
}}


\par }{\PP \VS{19}{\WJ{Cunosc faptele voastre, dragostea voastră, credința, slujirea, răbdarea și răbdarea, și știu că ultimele voastre fapte sunt mai mari decât cele dintâi.
}}

\VS{20}{\WJ{Dar am un lucru împotriva voastră: că o tolerați pe
}}\FTNT{§}{{\FR 2:20 }{\FT{TR, NU citiți “că” în loc de “dumneavoastră”}}}{\WJ{femeia voastră Izabela, care se numește proorociță. Ea îi învață și îi seduce pe slujitorii mei să comită imoralitate sexuală și să mănânce lucruri sacrificate idolilor.
}}

\VS{21}I-am
{\WJ{dat timp să se pocăiască, dar ea refuză să se pocăiască de imoralitatea ei sexuală.
}}

\VS{22}{\WJ{Iată că o voi arunca pe ea și pe cei care comit adulter cu ea într-un pat de mare asuprire, dacă nu se pocăiesc de faptele ei.
}}

\VS{23}{\WJ{Îi voi ucide copiii cu moartea și toate adunările vor ști că Eu sunt Cel care cercetează mințile și inimile. Voi da fiecăruia dintre voi după faptele sale.
}}

\VS{24}{\WJ{Dar vouă vă spun: celorlalți care sunt în Thyatira — celor care nu au această învățătură, care nu cunosc ceea ce unii numesc “lucrurile profunde ale Satanei” — vouă vă spun: nu pun nicio altă povară asupra voastră.
}}

\VS{25}{\WJ{Totuși, țineți cu tărie ceea ce aveți până când voi veni.
}}

\VS{26}Cel care va
{\WJ{birui și cel care va păstra lucrările Mele până la sfârșit, aceluia îi voi da autoritate asupra națiunilor.
}}

\VS{27}{\WJ{El le va stăpâni cu un toiag de fier, sfărâmându-le ca pe niște vase de lut,
}}\FTNT{*}{{\FR 2:27 }{\FT{{\CHAR{Psalmul 2:9}}}}}{\WJ{cum am primit și eu de la Tatăl Meu;
}}

\VS{28}{\WJ{și-i voi da steaua dimineții.
}}

\VS{29}{\WJ{Cine are urechi, să asculte ce spune Duhul Sfânt adunărilor.
}}



\par }\Chap{3}{\PP \VerseOne{1}{\WJ{“Și scrie îngerului adunării din Sardes:
}}

\par }{\PP {\WJ{“Cel care are cele șapte duhuri ale lui Dumnezeu și cele șapte stele spune aceste lucruri:
}}

\par }{\PP {\WJ{“Îți cunosc faptele, că ai reputația de a fi viu, dar ești mort.
}}

\VS{2}{\WJ{Trezește-te și întărește ceea ce ți-a rămas, pe care erai pe punctul de a lepăda,
}}\FTNT{*}{{\FR 3:2 }{\FT{NU \& TR citesc “care erau pe cale să moară” în loc de “pe care erați pe cale să le aruncați”.}}}{\WJ{căci nu am găsit nici o lucrare de-a ta desăvârșită înaintea Dumnezeului meu.
}}

\VS{3}{\WJ{Amintiți-vă, așadar, cum ați primit și cum ați auzit. Păstrați-o și pocăiți-vă. De aceea, dacă nu veți veghea, voi veni ca un hoț și nu veți ști la ce oră voi veni peste voi.
}}

\VS{4}Cu toate
{\WJ{acestea, aveți câteva nume în Sardes care nu și-au pângărit hainele. Ei vor umbla cu mine în alb, pentru că sunt vrednici.
}}

\VS{5}Cel care va
{\WJ{învinge va fi îmbrăcat în haine albe și nu-i voi șterge nicidecum numele din cartea vieții și-i voi mărturisi numele înaintea Tatălui Meu și înaintea îngerilor Lui.
}}

\VS{6}{\WJ{Cine are urechi, să asculte ce spune Duhul Sfânt adunărilor.
}}


\par }{\PP \VS{7}{\WJ{“Scrieți îngerului adunării din Filadelfia:
}}

\par }{\PP {\WJ{“Cel ce este sfânt, cel adevărat, cel ce are cheia lui David, cel ce deschide și nimeni nu poate închide, cel ce închide și nimeni nu deschide, spune aceste lucruri}}:


\par }{\PP \VS{8}“{\WJ{Știu faptele tale (iată, ți-am pus înaintea ta o ușă deschisă, pe care nimeni nu o poate închide), că ai puțină putere, că ai păzit cuvântul Meu și că nu te-ai lepădat de Numele Meu.
}}

\VS{9}{\WJ{Iată, voi face ca unii din sinagoga Satanei, dintre cei care spun că sunt iudei, și nu sunt, ci mint — iată, îi voi face să vină și să se închine înaintea picioarelor tale și să știe că te-am iubit.
}}

\VS{10}{\WJ{Pentru că ați păzit porunca mea de a răbda, vă voi păzi și eu de ceasul încercării care va veni peste întreaga lume, pentru a pune la încercare pe cei care locuiesc pe pământ.
}}

\VS{11}{\WJ{Eu vin repede! Țineți cu tărie ceea ce aveți, ca nimeni să nu vă ia coroana.
}}

\VS{12}Pe cel care va
{\WJ{birui, îl voi face stâlp în templul Dumnezeului Meu și nu va mai ieși de acolo. Voi scrie pe el numele Dumnezeului meu și numele cetății Dumnezeului meu, noul Ierusalim, care se coboară din cer de la Dumnezeul meu, și numele meu cel nou.
}}

\VS{13}{\WJ{Cine are urechi, să asculte ce spune Duhul Sfânt adunărilor.
}}


\par }{\PP \VS{14}{\WJ{“Scrieți îngerului adunării din Laodiceea:
}}

\par }{\PP {\WJ{“Amin, Martorul credincios și adevărat, Începutul}}\FTNT{†}{{\FR 3:14 }{\FT{sau, Sursa, sau Capul}}}{\WJ{creației lui Dumnezeu, spune următoarele:
}}


\par }{\PP \VS{15}{\WJ{Eu cunosc faptele tale, că nu ești nici rece, nici fierbinte. Aș vrea să fiți reci sau fierbinți.
}}

\VS{16}{\WJ{Deci, pentru că sunteți călduți, și nu sunteți nici calzi, nici reci, nici fierbinți, vă voi vomita din gura mea.
}}

\VS{17}{\WJ{Pentru că zici: “Sunt bogat, am dobândit bogății și nu am nevoie de nimic”, și nu știi
}}că{\WJ{tu ești un nenorocit, mizerabil, sărac, orb și gol;
}}

\VS{18}{\WJ{te sfătuiesc să cumperi de la mine aur rafinat prin foc, ca să te îmbogățești, și haine albe, ca să te îmbraci și să nu se vadă rușinea goliciunii tale, și unguent pentru ochii tăi, ca să poți vedea.
}}

\VS{19}Pe
{\WJ{toți cei pe care îi iubesc, îi mustru și îi pedepsesc. Fiți deci zeloși și pocăiți-vă.
}}

\VS{20}{\WJ{Iată, Eu stau la ușă și bat. Dacă aude cineva glasul Meu și deschide ușa, atunci voi intra la el și voi lua masa cu el, și el cu Mine.
}}

\VS{21}Celui ce va
{\WJ{birui, îi voi da să șadă cu Mine pe tronul Meu, după cum și Eu am biruit și am șezut cu Tatăl Meu pe tronul Lui.
}}

\VS{22}{\WJ{Cine are urechi, să asculte ce spune Duhul Sfânt adunărilor.”
}}



\par }\Chap{4}{\PP \VerseOne{1}După acestea, m-am uitat și am văzut o ușă deschisă în cer; și primul glas pe care l-am auzit, ca o trâmbiță care vorbea cu mine, a fost un glas care zicea: “Urcă-te aici și-ți voi arăta cele ce trebuie să se întâmple după aceasta.”


\par }{\PP \VS{2}Îndată am fost în Duhul Sfânt. Și iată că era un tron așezat în cer și pe tron ședea cineva

\VS{3}care semăna cu o piatră de jasp și cu un sardiu. În jurul tronului era un curcubeu, ca un smarald la care se putea privi.

\VS{4}În jurul tronului erau douăzeci și patru de scaune de domnie. Pe aceste tronuri ședeau douăzeci și patru de bătrâni, îmbrăcați în haine albe, cu coroane de aur pe cap.

\VS{5}Din tron ies fulgere, sunete și tunete. Înaintea tronului său ardeau șapte lămpi de foc, care sunt cele șapte duhuri ale lui Dumnezeu.

\VS{6}Înaintea tronului era ceva asemănător cu o mare de sticlă, asemănătoare cristalului. În mijlocul tronului și în jurul tronului erau patru ființe vii pline de ochi în față și în spate.

\VS{7}Prima făptură era ca un leu, a doua făptură ca un vițel, a treia făptură avea o față de om, iar a patra era ca un vultur zburător.

\VS{8}Cele patru făpturi vii, fiecare dintre ele având șase aripi, sunt pline de ochi în jur și înăuntru. Ele nu au odihnă nici ziua, nici noaptea, spunând: “Sfânt, sfânt, sfânt\FTNT{*}{{\FR 4:8 }{\FT{Hodges/Farstad MT citește “sfânt” de 9 ori în loc de 3.}}} este Domnul Dumnezeu, Atotputernicul, cel care a fost, cel care este și cel care va veni!”


\par }{\PP \VS{9}Când făpturile vii vor da slavă, cinste și mulțumire Celui ce șade pe tron, Celui ce trăiește în vecii vecilor,

\VS{10}cei douăzeci și patru de bătrâni se vor arunca înaintea Celui ce șade pe tron și se vor închina Celui ce trăiește în vecii vecilor și își vor arunca coroanele înaintea tronului, zicând:

\VS{11}“Vrednic ești Tu, Domnul și Dumnezeul nostru, Cel Sfânt,
\FTNT{†}{{\FR 4:11 }{\FT{TR omite “și Dumnezeu, Cel Sfânt,”}}}să primești slava, cinstea și puterea, căci Tu ai creat toate lucrurile și, din voia Ta, ele au existat și au fost create!”.



\par }\Chap{5}{\PP \VerseOne{1}Am văzut în mâna dreaptă a Celui ce ședea pe scaunul de domnie o carte scrisă pe dinăuntru și pe dinafară, închisă cu șapte peceți.

\VS{2}Am văzut un înger puternic care striga cu glas tare: “Cine este vrednic să deschidă cartea și să-i rupă pecețile?”

\VS{3}Nimeni din cerurile de sus, de pe pământ sau de sub pământ nu era în stare să deschidă cartea sau să se uite în ea.

\VS{4}Atunci am plâns mult, pentru că nimeni nu a fost găsit vrednic să deschidă cartea sau să se uite în ea.

\VS{5}Unul dintre bătrâni mi-a spus: “Nu plânge. Iată că a biruit Leul care este din seminția lui Iuda, Rădăcina lui David: cel care deschide cartea și cele șapte peceți ale ei.”


\par }{\PP \VS{6}Și am văzut în mijlocul tronului și al celor patru făpturi vii și în mijlocul bătrânilor, un Miel în picioare, ca și cum ar fi fost înjunghiat, având șapte coarne și șapte ochi, care sunt cele șapte duhuri ale lui Dumnezeu, trimise pe tot pământul.

\VS{7}Apoi a venit și a luat-o din mâna dreaptă a Celui care ședea pe tron.

\VS{8}După ce a luat cartea, cele patru făpturi vii și cei douăzeci și patru de bătrâni au căzut înaintea Mielului, fiecare având o harpă și cupe de aur pline de tămâie, care sunt rugăciunile sfinților.

\VS{9}Și au cântat o cântare nouă, zicând: “Nu, nu!

\par }{\Q “Ești vrednic să iei cartea

\par }{\QB și să-i deschidă pecețile,

\par }{\Q pentru că ai fost ucis,

\par }{\QB și ne-ai cumpărat pentru Dumnezeu cu sângele tău

\par }{\QB din orice seminție, limbă, popor și națiune,


\par }{\Q \VS{10}și ne-a făcut împărați și preoți pentru Dumnezeul nostru;

\par }{\QB și vom domni pe pământ.”


\par }{\PP \VS{11}M-am uitat și am auzit un glas ca un glas de mulți îngeri în jurul scaunului de domnie, al făpturilor vii și al bătrânilor. Numărul lor era de zece mii de zeci de mii și de mii de mii,

\VS{12}care spuneau cu glas tare: “Vrednic este Mielul care a fost ucis să primească puterea, bogăția, înțelepciunea, tăria, onoarea, gloria și binecuvântarea!”


\par }{\PP \VS{13}Și am auzit toate făpturile care sunt în ceruri, pe pământ, sub pământ, pe mare și tot ce este în ele, zicând: “Binecuvântarea, cinstea, slava și stăpânirea să fie pentru Cel ce șade pe scaunul de domnie și pentru Miel, în vecii vecilor! Amin!”
\FTNT{*}{{\FR 5:13 }{\FT{TR omite “Amin!”}}}


\par }{\PP \VS{14}Cele patru făpturi vii au zis: “Amin!” Apoi
\FTNT{†}{{\FR 5:14 }{\FT{TR adaugă “douăzeci și patru”}}}bătrânii au căzut la pământ și s-au închinat.
\FTNT{‡}{{\FR 5:14 }{\FT{TR adaugă “Cel ce trăiește în vecii vecilor”}}}



\par }\Chap{6}{\PP \VerseOne{1}Am văzut că Mielul a deschis una din cele șapte peceți și am auzit pe una din cele patru făpturi vii zicând, ca un glas de tunet: “Vino și vezi!”.

\VS{2}Apoi a apărut un cal alb, iar cel care ședea pe el avea un arc. I s-a dat o coroană și a ieșit învingător și să cucerească.


\par }{\PP \VS{3}Când a deschis pecetea a doua, am auzit a doua făptură vie care zicea: “Vino!”

\VS{4}A ieșit un alt cal roșu. Celui care ședea pe el i s-a dat puterea de a lua pacea de pe pământ și de a se ucide unii pe alții. I s-a dat o sabie mare.


\par }{\PP \VS{5}Când a deschis pecetea a treia, am auzit a treia făptură vie care zicea: “Vino și vezi!” Și iată un cal negru, iar cel care ședea pe el avea o balanță în mână.

\VS{6}Am auzit o voce în mijlocul celor patru făpturi vii care spunea: “Un choenix
\FTNT{*}{{\FR 6:6 }{\FT{Un choenix este o măsură de volum uscat care este puțin mai mare decât un litru (puțin mai mare decât un sfert de litru).}}}de grâu pentru un denar și trei choenix de orz pentru un denar! Nu stricați uleiul și vinul!”


\par }{\PP \VS{7}Când a deschis pecetea a patra, am auzit a patra făptură vie care zicea: “Vino și vezi!”

\VS{8}Și iată un cal palid, iar numele celui care ședea pe el era Moartea. Hades\FTNT{†}{{\FR 6:8 }{\FT{sau, iad}}} îl urma cu el. I s-a dat autoritate asupra unui sfert din pământ, ca să ucidă cu sabia, cu foametea, cu moartea și cu animalele sălbatice de pe pământ.


\par }{\PP \VS{9}Când a deschis pecetea a cincea, am văzut sub altar sufletele celor ce fuseseră uciși pentru Cuvântul lui Dumnezeu și pentru mărturia Mielului, pe care o aveau.

\VS{10}Ei strigau cu glas tare, zicând: “Până când, Stăpâne, Cel sfânt și adevărat, până când vei judeca și vei răzbuna sângele nostru pe cei care locuiesc pe pământ?”

\VS{11}O haină albă și lungă a fost dată fiecăruia dintre ei. Li s-a spus că ar trebui să se odihnească încă o vreme, până când tovarășii lor de slujbă și frații lor, care vor fi și ei uciși la fel ca și ei, își vor încheia cursul.


\par }{\PP \VS{12}Am văzut când a deschis pecetea a șasea și a fost un mare cutremur de pământ. Soarele s-a făcut negru ca un sac de păr și toată luna s-a făcut ca sângele.

\VS{13}Stelele cerului au căzut pe pământ, ca un smochin care-și lasă smochinele necoapte când este scuturat de un vânt mare.

\VS{14}Cerul a fost îndepărtat ca un sul de hârtie când este rulat. Fiecare munte și fiecare insulă au fost mutate de la locul lor.

\VS{15}Împărații pământului, prinții, comandanții, cei bogați, cei puternici și toți sclavii și oamenii liberi s-au ascuns în peșteri și în stâncile munților.

\VS{16}Ei au spus munților și stâncilor: “Cădeți peste noi și ascundeți-ne de fața Celui care șade pe tron și de mânia Mielului,

\VS{17}pentru că a venit ziua cea mare a mâniei Lui și cine poate sta în picioare?”



\par }\Chap{7}{\PP \VerseOne{1}După aceea am văzut patru îngeri care stăteau în picioare la cele patru colțuri ale pământului și care țineau cele patru vânturi ale pământului, ca să nu mai sufle nici un vânt pe pământ, nici pe mare, nici pe vreun copac.

\VS{2}Am văzut un alt înger care se înălța de la răsăritul soarelui, având pecetea Dumnezeului celui viu. El a strigat cu glas tare către cei patru îngeri cărora le fusese dat să facă rău pământului și mării,

\VS{3}spunând: “Nu faceți rău nici pământului, nici mării, nici copacilor, până când nu le vom pune pe frunțile lor sigiliul robilor Dumnezeului nostru!”

\VS{4}Am auzit numărul celor care au fost pecetluiți: o sută patruzeci și patru de mii, pecetluiți din fiecare seminție a copiilor lui Israel:


\par }{\Q \VS{5}Din seminția lui Iuda, douăsprezece mii au fost pecetluiți,

\par }{\Q din seminția lui Ruben, douăsprezece mii,

\par }{\Q din seminția lui Gad, douăsprezece mii,


\par }{\Q \VS{6}din seminția lui Așer douăsprezece mii,

\par }{\Q din seminția lui Neftali, douăsprezece mii,

\par }{\Q din seminția lui Manase, douăsprezece mii,


\par }{\Q \VS{7}din seminția lui Simeon, douăsprezece mii,

\par }{\Q din seminția lui Levi, douăsprezece mii,

\par }{\Q din seminția lui Isahar, douăsprezece mii,


\par }{\Q \VS{8}din seminția lui Zabulon douăsprezece mii,

\par }{\Q din seminția lui Iosif, douăsprezece mii, și

\par }{\Q din seminția lui Beniamin, douăsprezece mii au fost pecetluiți.


\par }{\PP \VS{9}După acestea, m-am uitat și iată o mulțime mare, pe care nimeni nu o putea număra, din toate neamurile și din toate semințiile, popoarele și limbile, care stătea în picioare înaintea scaunului de domnie și înaintea Mielului, îmbrăcați în haine albe, cu ramuri de palmier în mâini.

\VS{10}Ei strigau cu glas tare, zicând: “Mântuirea fie a Dumnezeului nostru, care șade pe tron, și a Mielului!”


\par }{\PP \VS{11}Toți îngerii stăteau în picioare în jurul scaunului de domnie, bătrânii și cele patru făpturi vii; și, căzând cu fața la pământ înaintea scaunului de domnie, s-au închinat lui Dumnezeu,

\VS{12}zicând: “Amin! Binecuvântare, glorie, înțelepciune, mulțumire, onoare, putere și tărie să fie pentru Dumnezeul nostru în vecii vecilor! Amin.”


\par }{\PP \VS{13}Unul din bătrâni a luat cuvântul și mi-a zis: “Aceștia care sunt îmbrăcați în haine albe, cine sunt și de unde au venit?”


\par }{\PP \VS{14}I-am spus: “Stăpâne, tu știi.”

\par }{\PP El mi-a spus: “Aceștia sunt cei care au ieșit din marea suferință.\FTNT{*}{{\FR 7:14 }{\FT{sau, asuprire}}} Ei și-au spălat veșmintele și le-au albit în sângele Mielului.

\VS{15}De aceea, ei sunt înaintea tronului lui Dumnezeu și îi slujesc zi și noapte în templul lui. Cel care șade pe tron își va întinde cortul peste ei.

\VS{16}Ei nu vor mai fi niciodată flămânzi sau însetați. Soarele nu va bate peste ei, nici căldura,

\VS{17}căci Mielul care este în mijlocul tronului îi paște și îi conduce la izvoare de ape dătătoare de viață. Și Dumnezeu va șterge orice lacrimă din ochii lor.”



\par }\Chap{8}{\PP \VerseOne{1}Când a deschis pecetea a șaptea, a fost liniște în cer timp de o jumătate de oră.

\VS{2}Am văzut cei șapte îngeri care stau înaintea lui Dumnezeu și li s-au dat șapte trâmbițe.


\par }{\PP \VS{3}Un alt înger a venit și a stat deasupra altarului, având un cădelniță de aur. I s-a dat multă tămâie, ca să o adauge la rugăciunile tuturor sfinților pe altarul de aur care era înaintea tronului.

\VS{4}Fumul tămâiei, împreună cu rugăciunile sfinților, se ridica înaintea lui Dumnezeu din mâna îngerului.

\VS{5}Îngerul a luat cădelnița, a umplut-o cu focul altarului, apoi a aruncat-o pe pământ. Au urmat tunete, sunete, fulgere și un cutremur.


\par }{\PP \VS{6}Cei șapte îngeri care aveau cele șapte trâmbițe s-au pregătit să sune.


\par }{\PP \VS{7}A sunat primul sunet și a urmat grindină și foc amestecat cu sânge, și au fost aruncați la pământ. O treime din pământ a ars,\FTNT{*}{{\FR 8:7 }{\FT{TR omite “O treime din pământ a fost arsă”}}} o treime din copaci au ars și toată iarba verde a ars.


\par }{\PP \VS{8}Al doilea înger a sunat din glas, și ceva ca un munte mare și arzător a fost aruncat în mare. O treime din mare s-a transformat în sânge,

\VS{9}și o treime din viețuitoarele care erau în mare au murit. O treime din corăbii au fost distruse.


\par }{\PP \VS{10}Îngerul al treilea a sunat și a căzut din cer o stea mare, care ardea ca o torță, și a căzut pe o treime din râuri și pe izvoarele de apă.

\VS{11}Numele stelei este “Viermele”. O treime din ape au devenit pelin de vierme. Mulți oameni au murit din cauza apelor, pentru că acestea au devenit amare.


\par }{\PP \VS{12}Îngerul al patrulea a sunat și a lovit o treime din soare, o treime din lună și o treime din stele, astfel încât o treime din ele să se întunece, iar ziua să nu mai strălucească o treime din ea și noaptea la fel.

\VS{13}Am văzut și am auzit un vultur care
\FTNT{†}{{\FR 8:13 }{\FT{TR citește “înger” în loc de “vultur”
}}}zbura prin mijlocul cerului și care spunea cu glas tare: “Vai! Vai! Vai de cei care locuiesc pe pământ, din cauza celorlalte sunete ale trâmbițelor celor trei îngeri, care urmează să sune!”



\par }\Chap{9}{\PP \VerseOne{1}Îngerul al cincilea a sunat din glas și am văzut din cer o stea care a căzut pe pământ. Cheia de la groapa abisului i-a fost dată.

\VS{2}El a deschis groapa abisului și din groapă a ieșit fum ca fumul unui cuptor
\FTNT{*}{{\FR 9:2 }{\FT{TR adaugă “mare”}}}aprins. Soarele și aerul s-au întunecat din cauza fumului care ieșea din groapă.

\VS{3}Din fum au ieșit lăcuste pe pământ și li s-a dat putere, cum au puterea scorpionilor de pe pământ.

\VS{4}Li s-a spus că nu trebuie să facă rău nici ierbii pământului, nici unei plante verzi, nici unui copac, ci numai acelor oameni care nu au pecetea lui Dumnezeu pe frunte.

\VS{5}Li s-a dat putere, nu pentru a-i ucide, ci pentru a-i chinui timp de cinci luni. Chinul lor a fost ca chinul unui scorpion atunci când lovește o persoană.

\VS{6}În zilele acelea, oamenii vor căuta moartea și nu o vor găsi în niciun fel. Ei vor dori să moară, iar moartea va fugi de ei.


\par }{\PP \VS{7}Lăcustele erau ca niște cai pregătiți pentru război. Pe capetele lor erau ceva ca niște coroane de aur, iar fețele lor erau ca fețele oamenilor.

\VS{8}Aveau părul ca părul femeilor, iar dinții lor erau ca cei ai leilor.

\VS{9}Aveau platoșe ca niște platoșe de fier. Zgomotul aripilor lor era ca zgomotul multor care și cai care se grăbeau la război.

\VS{10}Aveau cozi ca cele ale scorpionilor, cu înțepături. În cozile lor au puterea de a face rău oamenilor timp de cinci luni.

\VS{11}Ei au peste ei ca rege pe îngerul abisului. Numele lui în ebraică este “Abaddon”,
\FTNT{†}{{\FR 9:11 }{\FT{“Abaddon” este un cuvânt ebraic care înseamnă “ruină”, “distrugere” sau “locul distrugerii”.}}}dar în greacă are numele “Apollyon”.
\FTNT{‡}{{\FR 9:11 }{\FT{“Apollyon” înseamnă “Distrugător”.}}}


\par }{\PP \VS{12}Primul necaz a trecut. Iată, mai sunt încă două nenorociri care urmează după aceasta.


\par }{\PP \VS{13}Îngerul al șaselea a sunat. Am auzit un glas din coarnele altarului de aur care este înaintea lui Dumnezeu,

\VS{14}care spunea celui de-al șaselea înger care avea trâmbița: “Eliberați pe cei patru îngeri care sunt legați la râul cel mare Eufrat!”


\par }{\PP \VS{15}Au fost eliberați cei patru îngeri care fuseseră pregătiți pentru ceasul, ziua, luna și anul acela, ca să ucidă o treime din omenire.

\VS{16}Numărul armatelor călăreților era de două sute de milioane.\FTNT{§}{{\FR 9:16 }{\FT{literal, “zece mii de zece mii”}}} Am auzit numărul lor.

\VS{17}Astfel am văzut în viziune caii și pe cei care ședeau pe ei, având platoșe de un roșu aprins, de un albastru de iacint și de un galben de sulf; iar capetele cailor semănau cu capetele leilor. Din gurile lor ieșea foc, fum și sulf.

\VS{18}Prin aceste trei plăgi, o treime din omenire a fost ucisă: de focul, de fumul și de sulf, care ieșeau din gurile lor.

\VS{19}Căci puterea cailor este în gurile și în cozile lor. Căci cozile lor sunt ca niște șerpi și au capete, și cu ele fac rău.


\par }{\PP \VS{20}Ceilalți oameni, care nu au fost uciși de aceste plăgi, nu s-au pocăit de faptele mâinilor lor, ca să nu se închine demonilor și idolilor de aur, de argint, de aramă, de piatră și de lemn, care nu văd, nu aud și nu umblă.

\VS{21}Nu s-au pocăit de uciderile lor, de vrăjitoriile lor,\FTNT{*}{{\FR 9:21 }{\FT{Cuvântul pentru “vrăjitorii” (pharmakeia) implică, de asemenea, folosirea de poțiuni, otrăvuri și medicamente.}}} de imoralitatea lor sexuală și de furturile lor.



\par }\Chap{10}{\PP \VerseOne{1}Am văzut un înger puternic coborând din cer, îmbrăcat într-un nor. Un curcubeu era pe capul lui. Fața lui era ca soarele și picioarele lui erau ca niște stâlpi de foc.

\VS{2}El avea în mână o carte mică și deschisă. Și-a pus piciorul drept pe mare și piciorul stâng pe pământ.

\VS{3}A strigat cu glas tare, cum răcnește un leu. Când a strigat, cele șapte tunete și-au scos glasurile.

\VS{4}Când au răsunat cele șapte tunete, eram pe punctul de a scrie, dar am auzit un glas din cer care spunea: “Sigilează lucrurile pe care le-au spus cele șapte tunete și nu le scrie.”


\par }{\PP \VS{5}Îngerul pe care l-am văzut stând în picioare pe mare și pe uscat și-a ridicat mâna dreaptă spre cer

\VS{6}și a jurat pe Cel ce trăiește în vecii vecilor, care a creat cerul și lucrurile care sunt în el, pământul și lucrurile care sunt în el, marea și lucrurile care sunt în ea, că nu va mai fi nici o întârziere,

\VS{7}ci în zilele glasului celui de-al șaptelea înger, când va suna, atunci se va sfârși taina lui Dumnezeu, după cum a spus robilor Săi, proorocii.


\par }{\PP \VS{8}Glasul pe care l-am auzit din cer, care vorbea din nou cu mine, a zis: “Du-te și ia cartea deschisă în mâna îngerului care stă pe mare și pe uscat.”


\par }{\PP \VS{9}M-am dus la înger și i-am spus să-mi dea cărticica.

\par }{\PP Mi-a spus: “Ia-l și mănâncă. Îți va amărî stomacul, dar în gură va fi dulce ca mierea”.


\par }{\PP \VS{10}Am luat cărticica din mâna îngerului și am mâncat-o. A fost dulce ca mierea în gura mea. După ce am mâncat-o, stomacul mi s-a făcut amar.

\VS{11}Ei\FTNT{*}{{\FR 10:11 }{\FT{TR citește “El” în loc de “Ei”}}} mi-au spus: “Trebuie să profețești din nou peste multe popoare, națiuni, limbi și împărați.”



\par }\Chap{11}{\PP \VerseOne{1}Mi-a fost dată o trestie ca o vergea. Cineva a zis: “Ridică-te și măsoară templul lui Dumnezeu, altarul și pe cei care se închină în el.

\VS{2}Lasă deoparte curtea care este în afara templului și nu o măsura, pentru că ea a fost dată națiunilor. Ele vor călca în picioare cetatea sfântă timp de patruzeci și două de luni.

\VS{3}Voi da putere celor doi martori ai Mei și ei vor profeți o mie două sute șaizeci de zile, îmbrăcați în sac.”


\par }{\PP \VS{4}Aceștia sunt cei doi măslini și cele două sfeșnice, care stau înaintea Domnului pământului.

\VS{5}Dacă cineva vrea să le facă rău, focul iese din gura lor și-i mistuie pe vrăjmașii lor. Dacă cineva vrea să le facă rău, trebuie să fie ucis în felul acesta.

\VS{6}Aceștia au puterea să închidă cerul, ca să nu plouă în zilele profeției lor. Ei au putere asupra apelor, ca să le transforme în sânge și să lovească pământul cu orice fel de molimă, de câte ori vor ei.


\par }{\PP \VS{7}După ce-și vor fi isprăvit mărturia lor, fiara care se va ridica din prăpastie va face război cu ei, îi va birui și-i va ucide.

\VS{8}Trupurile lor moarte vor fi pe strada marii cetăți, care, din punct de vedere spiritual, se numește Sodoma și Egipt, unde a fost răstignit și Domnul lor.

\VS{9}Dintre popoare, triburi, limbi și națiuni, oamenii se vor uita la cadavrele lor timp de trei zile și jumătate și nu vor îngădui să li se pună cadavrele într-un mormânt.

\VS{10}Cei care locuiesc pe pământ se vor bucura pentru ei și se vor veseli. Își vor face daruri unii altora, pentru că acești doi profeți i-au chinuit pe cei care locuiesc pe pământ.


\par }{\PP \VS{11}După cele trei zile și jumătate, a intrat în ei suflare de viață de la Dumnezeu și s-au ridicat în picioare. O mare frică a cuprins pe cei care i-au văzut.

\VS{12}Am auzit un glas puternic din cer care le spunea: “Veniți aici sus!” Ei s-au suit la cer într-un nor și dușmanii lor i-au văzut.

\VS{13}În ziua aceea a fost un mare cutremur și o zecime din cetate a căzut. Șapte mii de oameni au fost uciși în cutremur, iar restul s-au îngrozit și au dat slavă Dumnezeului cerului.


\par }{\PP \VS{14}A doua nenorocire a trecut. Iată, al treilea vai vine repede.


\par }{\PP \VS{15}Îngerul al șaptelea a sunat și au urmat glasuri mari în ceruri, care ziceau: “Împărăția lumii a devenit Împărăția Domnului nostru și a Hristosului Său. El va domni în vecii vecilor!”


\par }{\PP \VS{16}Cei douăzeci și patru de bătrâni, care stau pe tronurile lor înaintea scaunului de domnie al lui Dumnezeu, s-au aruncat cu fața la pământ și s-au închinat lui Dumnezeu,

\VS{17}zicând: “Îți mulțumim, Doamne Dumnezeule, Cel Atotputernic, Cel care este și care era,
\FTNT{*}{{\FR 11:17 }{\FT{TR adaugă “și care vine”}}}pentru că ți-ai luat marea ta putere și ai domnit.

\VS{18}Neamurile s-au mâniat, iar mânia ta a venit, ca și timpul ca morții să fie judecați și să le dai robilor tăi, profeții, răsplata lor, precum și sfinților și celor care se tem de numele tău, celor mici și celor mari, și să nimicești pe cei care distrug pământul.”


\par }{\PP \VS{19}Templul lui Dumnezeu, care este în ceruri, s-a deschis și chivotul legământului Domnului s-a văzut în templul lui. Au urmat fulgere, zgomote, tunete, cutremur și grindină mare.



\par }\Chap{12}{\PP \VerseOne{1}Un semn mare s-a văzut în cer: o femeie îmbrăcată cu soarele și cu luna sub picioarele ei, iar pe cap o cunună de douăsprezece stele.

\VS{2}Ea era însărcinată. Ea a strigat de durere, chinuindu-se să nască.


\par }{\PP \VS{3}Un alt semn s-a văzut în cer. Iată un balaur mare și roșu, care avea șapte capete și zece coarne, și pe capetele lui șapte coroane.

\VS{4}Coada lui a tras o treime din stelele cerului și le-a aruncat pe pământ. Balaurul stătea în fața femeii care urma să nască, pentru ca, atunci când aceasta va naște, să-i devoreze copilul.

\VS{5}Ea a născut un fiu, un copil de parte bărbătească, care va conduce toate națiunile cu un toiag de fier. Copilul ei a fost răpit la Dumnezeu și la tronul Lui.

\VS{6}Femeia a fugit în pustiu, unde are un loc pregătit de Dumnezeu, pentru ca acolo să o hrănească o mie două sute șaizeci de zile.


\par }{\PP \VS{7}Pe cer era război. Mihail și îngerii săi au făcut război cu balaurul. Balaurul și îngerii lui au făcut război.

\VS{8}Ei nu au biruit. Nu s-a mai găsit loc pentru ei în cer.

\VS{9}A fost doborât balaurul cel mare, șarpele cel vechi, cel care se numește diavolul și Satana, înșelătorul întregii lumi. El a fost aruncat pe pământ, iar îngerii lui au fost aruncați împreună cu el.


\par }{\PP \VS{10}Și am auzit în cer un glas tare, care zicea: “Acum a venit mântuirea, puterea și Împărăția Dumnezeului nostru și autoritatea Hristosului Său, căci a fost doborât acuzatorul fraților noștri, care-i acuza zi și noapte înaintea Dumnezeului nostru.

\VS{11}Ei l-au biruit datorită sângelui Mielului și datorită cuvântului mărturiei lor. Ei nu și-au iubit viața, până la moarte.

\VS{12}De aceea, bucurați-vă, ceruri și voi, care locuiți în ele! Vai de pământ și de mare, pentru că diavolul s-a coborât la voi, având o mare mânie, știind că nu are decât puțin timp.”


\par }{\PP \VS{13}Când a văzut balaurul că a fost aruncat pe pământ, a prigonit pe femeia care a născut copilul de parte bărbătească.

\VS{14}Femeii i-au fost date două aripi de vultur mare, ca să zboare în pustiu, la locul ei, ca să se hrănească o vreme, vremuri și jumătate de vreme, departe de fața șarpelui.

\VS{15}Șarpele a scuipat apă din gură după femeie, ca un râu, ca să o facă să fie luată de curent.

\VS{16}Pământul a ajutat-o pe femeie, iar pământul și-a deschis gura și a înghițit râul pe care balaurul îl arunca din gura lui.

\VS{17}Balaurul s-a mâniat pe femeie și a plecat să se războiască cu restul urmașilor ei,
\FTNT{*}{{\FR 12:17 }{\FT{sau, sămânță}}}care păzesc poruncile lui Dumnezeu și țin mărturia lui Isus.



\par }\Chap{13}{\PP \VerseOne{1}Atunci am stat pe nisipul mării. Am văzut o fiară care se ridica din mare, având zece coarne și șapte capete. Pe coarnele sale erau zece coroane, iar pe capetele sale, nume blasfemiatoare.

\VS{2}Fiara pe care am văzut-o semăna cu un leopard, picioarele îi erau ca ale unui urs, iar gura îi era ca gura unui leu. Balaurul i-a dat puterea, tronul și o mare autoritate.

\VS{3}Unul dintre capetele sale părea că fusese rănit mortal. Rana lui fatală a fost vindecată și tot pământul s-a mirat de fiară.

\VS{4}Ei se închinau balaurului, pentru că el își dăduse autoritatea fiarei; și se închinau fiarei, spunând: “Cine este ca fiara? Cine este în stare să facă război cu ea?”


\par }{\PP \VS{5}I-a fost dată o gură care vorbește lucruri mari și blasfemii. I s-a dat puterea de a face război timp de patruzeci și două de luni.

\VS{6}Și-a deschis gura pentru a blasfemia împotriva lui Dumnezeu, pentru a huli numele Lui, locuința Lui și pe cei ce locuiesc în ceruri.

\VS{7}I-a fost dat să facă război cu sfinții și să îi învingă. I s-a dat autoritate asupra oricărui trib, popor, limbă și națiune.

\VS{8}Toți cei care locuiesc pe pământ i se vor închina, toți cei al căror nume nu a fost scris de la întemeierea lumii în cartea vieții Mielului care a fost ucis.

\VS{9}Dacă are cineva urechi, să audă.

\VS{10}Dacă cineva trebuie să meargă în captivitate, acela va merge în captivitate. Dacă cineva trebuie să fie ucis cu sabia, trebuie să fie ucis.\FTNT{*}{{\FR 13:10 }{\FT{TR citește: “Dacă cineva duce în captivitate, în captivitate se duce. Dacă cineva va ucide cu sabia, trebuie să fie ucis cu sabia.” în loc de “Dacă cineva va merge în captivitate, va merge în captivitate. Dacă cineva va fi ucis cu sabia, trebuie să fie ucis”.
}}} Iată rezistența și credința sfinților.


\par }{\PP \VS{11}Am văzut o altă fiară care se ridica de pe pământ. Avea două coarne ca un miel și vorbea ca un balaur.

\VS{12}Ea exercita toată autoritatea primei fiare în prezența ei. Ea face ca pământul și cei care locuiesc pe el să se închine primei fiare, a cărei rană mortală a fost vindecată.

\VS{13}El face semne mari, făcând chiar să coboare foc din cer pe pământ în văzul oamenilor.

\VS{14}El înșală pe
\FTNT{†}{{\FR 13:14 }{\FT{NU omite “propria mea”}}}poporul meu care locuiește pe pământ, din cauza semnelor pe care i s-a îngăduit să le facă în fața fiarei, spunând celor care locuiesc pe pământ să facă un chip fiarei care avea rana de sabie și care a trăit.

\VS{15}I-a fost dat să dea suflare chipului fiarei, ca chipul fiarei să vorbească și să facă să fie uciși toți cei care nu se vor închina chipului fiarei.

\VS{16}El face ca tuturor, celor mici și celor mari, celor bogați și celor săraci, celor liberi și celor sclavi, să li se dea semne pe mâna dreaptă sau pe frunte;

\VS{17}și nimeni nu va putea să cumpere sau să vândă dacă nu are acel semn, care este numele fiarei sau numărul numelui ei.

\VS{18}Iată înțelepciunea. Cine are pricepere, să calculeze numărul fiarei, căci este numărul unui om. Numărul ei este șase sute șaizeci și șase.



\par }\Chap{14}{\PP \VerseOne{1}Am văzut și iată Mielul stând în picioare pe muntele Sionului, și împreună cu El un număr de o sută patruzeci și patru de mii, care aveau numele Lui și numele Tatălui Său scris pe frunțile lor.

\VS{2}Am auzit din cer un sunet ca un zgomot de ape multe și ca un tunet mare. Sunetul pe care l-am auzit era ca al harpiștilor care cântă la harpele lor.

\VS{3}Ei cântă o cântare nouă înaintea tronului, înaintea celor patru ființe vii și a bătrânilor. Nimeni nu putea să învețe cântecul, în afară de cei o sută patruzeci și patru de mii, cei care fuseseră răscumpărați de pe pământ.

\VS{4}Aceștia sunt cei care nu au fost spurcați cu femei, căci sunt fecioare. Aceștia sunt cei care îl urmează pe Miel oriunde merge. Aceștia au fost răscumpărați de Isus dintre oameni, cele dintâi roade pentru Dumnezeu și pentru Miel.

\VS{5}În gura lor nu s-a găsit nici o minciună, căci sunt ireproșabili.
\FTNT{*}{{\FR 14:5 }{\FT{TR adaugă “înaintea tronului lui Dumnezeu”}}}


\par }{\PP \VS{6}Și am văzut un înger care zbura din mijlocul cerului și care avea o veste veșnică, pe care trebuia să o vestească celor ce locuiesc pe pământ, oricărui neam, seminție, limbă și popor.

\VS{7}El a spus cu glas tare: “Temeți-vă de Domnul și dați-i slavă, căci a venit ceasul judecății Lui. Adorați-l pe cel care a făcut cerul, pământul, marea și izvoarele de apă!”.


\par }{\PP \VS{8}Și a urmat un al doilea înger, care a zis: “A căzut Babilonul cel mare, care a făcut toate neamurile să bea din vinul mâniei desfrânării lui.”


\par }{\PP \VS{9}Un alt înger, al treilea, i-a urmat și a zis cu glas mare: “Dacă cineva se închină fiarei și chipului ei și primește un semn pe frunte sau pe mână,

\VS{10}va bea și el din vinul mâniei lui Dumnezeu, care este pregătit neamestecat în paharul mâniei Lui. El va fi chinuit cu foc și sulf în prezența sfinților îngeri și în prezența Mielului.

\VS{11}Fumul chinului lor se va ridica în vecii vecilor. Nu au odihnă zi și noapte, cei care se închină fiarei și chipului ei și oricine primește semnul numelui ei.


\par }{\PP \VS{12}Aici este perseverența sfinților, a celor ce păzesc poruncile lui Dumnezeu și credința lui Isus.”


\par }{\PP \VS{13}Și am auzit o voce din ceruri care zicea: “Scrie: “Fericiți cei ce mor în Domnul de acum înainte!””.

\par }{\PP “Da”, spune Duhul Sfânt, “ca să se odihnească de munca lor, căci lucrările lor îi însoțesc”.


\par }{\PP \VS{14}M-am uitat și am văzut un nor alb și, pe nor, un om șezând ca un fiu de om,\FTNT{†}{{\FR 14:14 }{\FT{Daniel 7:13}}} având pe cap o cunună de aur și în mână o seceră ascuțită.

\VS{15}Un alt înger a ieșit din templu și a strigat cu glas tare către cel care ședea pe nor: “Trimite secera și seceră, căci a venit ceasul secerișului, căci secerișul pământului este copt!”

\VS{16}Cel care ședea pe nor a împins secera în pământ, și pământul a fost secerat.


\par }{\PP \VS{17}Un alt înger a ieșit din templul care este în cer. Avea și el o seceră ascuțită.

\VS{18}Un alt înger a ieșit din altar, cel care are putere asupra focului, și a strigat cu glas mare către cel care avea secera ascuțită, zicând: “Trimite secera ta ascuțită și culege ciorchinii viței de vie de pe pământ, căci strugurii pământului sunt pe deplin coapte!”

\VS{19}Îngerul și-a înfipt secera în pământ, a cules strugurii de pe pământ și i-a aruncat în marea presă de vin a mâniei lui Dumnezeu.

\VS{20}Presa de vin a fost călcată în afara cetății, și sângele a ieșit din ea până la frâiele cailor, până la o mie șase sute de stadii.
\FTNT{‡}{{\FR 14:20 }{\FT{1600 de stadii = 296 kilometri sau 184 mile}}}



\par }\Chap{15}{\PP \VerseOne{1}Am văzut pe cer un alt semn mare și minunat: șapte îngeri care aveau cele șapte plăgi de pe urmă, căci în ei s-a sfârșit mânia lui Dumnezeu.


\par }{\PP \VS{2}Am văzut ceva ca o mare de sticlă amestecată cu foc și pe cei ce au biruit fiara, chipul ei
\FTNT{*}{{\FR 15:2 }{\FT{TR adaugă “semnul său”.}}}și numărul numelui ei, care stăteau pe marea de sticlă, având harpele lui Dumnezeu.

\VS{3}Ei cântau cântecul lui Moise, robul lui Dumnezeu, și cântecul Mielului, zicând

\par }{\Q “Mari și minunate sunt lucrările Tale, Doamne Dumnezeule, Cel Atotputernic!

\par }{\QB Drepte și adevărate sunt căile Tale, Împărate al neamurilor.


\par }{\Q \VS{4}Cine nu s-ar teme de Tine, Doamne?

\par }{\QB și să-ți glorificăm numele?

\par }{\Q Căci numai Tu ești sfânt.

\par }{\QB Căci toate neamurile vor veni și se vor închina înaintea ta.

\par }{\QB Căci faptele Tale drepte au fost descoperite.”


\par }{\PP \VS{5}După acestea, m-am uitat, și templul cortului mărturiei din ceruri s-a deschis.

\VS{6}Cei șapte îngeri care aveau cele șapte plăgi au ieșit, îmbrăcați în in pur și strălucitor și purtând fâșii de aur în jurul pieptului.


\par }{\PP \VS{7}Una din cele patru făpturi vii a dat celor șapte îngeri șapte cupe de aur pline cu mânia lui Dumnezeu, care trăiește în vecii vecilor.

\VS{8}Templul s-a umplut de fum de la gloria lui Dumnezeu și de la puterea Lui. Nimeni nu a putut intra în templu până când cele șapte plăgi ale celor șapte îngeri nu se vor fi terminat.



\par }\Chap{16}{\PP \VerseOne{1}Am auzit un glas puternic din templu, care zicea celor șapte îngeri: “Duceți-vă și vărsați pe pământ cele șapte potire ale mâniei lui Dumnezeu!”.


\par }{\PP \VS{2}Cel dintâi s-a dus și a vărsat cupa lui pe pământ, care a devenit o rană dureroasă și neplăcută pentru oamenii care aveau semnul fiarei și care se închinau chipului ei.


\par }{\PP \VS{3}Al doilea înger a vărsat cupa sa în mare, și s-a făcut sânge ca de mort. Orice ființă vie din mare a murit.


\par }{\PP \VS{4}Al treilea a vărsat cupa sa în râuri și în izvoare de apă, și acestea s-au făcut sânge.

\VS{5}Am auzit pe îngerul apelor spunând: “Drept ești Tu, care ești și care erai, Sfinte, pentru că ai judecat aceste lucruri.

\VS{6}Căci ei au vărsat sângele sfinților și al profeților, iar tu le-ai dat să bea sânge. Ei merită acest lucru”.


\par }{\PP \VS{7}Am auzit altarul zicând: “Da, Doamne Dumnezeule, Atotputernicul, adevărate și drepte sunt judecățile Tale!”.


\par }{\PP \VS{8}Al patrulea și-a vărsat cupa pe soare și i s-a dat să ardă pe oameni cu foc.

\VS{9}Oamenii au fost arși de o căldură mare, iar oamenii au hulit numele lui Dumnezeu, care are putere asupra acestor plăgi. Ei nu s-au pocăit și nu i-au dat slavă.


\par }{\PP \VS{10}Al cincilea și-a vărsat cupa pe scaunul de domnie al fiarei, și împărăția ei s-a întunecat. Și-au ros limbile din cauza durerilor,

\VS{11}și au hulit pe Dumnezeul cerului din cauza durerilor și a rănilor lor. Și totuși nu s-au pocăit de faptele lor.


\par }{\PP \VS{12}Al șaselea și-a vărsat cupa pe râul cel mare, Eufrat. A secat apele lui, ca să fie pregătită calea pentru împărații care vor veni de la răsăritul soarelui.
\FTNT{*}{{\FR 16:12 }{\FT{sau, est}}}

\VS{13}Am văzut ieșind din gura balaurului, din gura fiarei și din gura falsului profet, trei duhuri necurate, ceva asemănător cu broaștele;

\VS{14}pentru că sunt duhuri de demoni, care fac semne și care merg la regii întregului pământ locuit, ca să-i adune pentru războiul din ziua aceea mare a lui Dumnezeu cel Atotputernic.


\par }{\PP \VS{15}“{\WJ{Iată, vin ca un hoț. Ferice de cel care veghează și își păstrează hainele, ca să nu umble gol și să nu i se vadă rușinea.”
}}

\VS{16}El i-a adunat în locul care se numește în ebraică “Harmaghedon”.


\par }{\PP \VS{17}Al șaptelea și-a vărsat cupa în aer. Un glas puternic a ieșit din templul cerului, de pe tron, spunând: “S-a făcut!”.

\VS{18}Au fost fulgere, sunete și tunete; și a fost un cutremur mare, cum nu s-a mai întâmplat de când există oameni pe pământ — un cutremur atât de mare și atât de puternic.

\VS{19}Marele oraș a fost împărțit în trei părți și cetățile națiunilor au căzut. Babilonul cel mare a fost amintit în ochii lui Dumnezeu, ca să-i dea paharul cu vinul de foc al mâniei Lui.

\VS{20}Toate insulele au fugit, și munții nu au mai fost găsiți.

\VS{21}Grindină mare, de greutatea unui talant,\FTNT{†}{{\FR 16:21 }{\FT{Un talent reprezintă aproximativ 30 de kilograme sau 66 de lire sterline.}}} a căzut din cer peste oameni. Oamenii l-au hulit pe Dumnezeu din cauza plăgii grindinei, căci această plagă a fost extrem de gravă.



\par }\Chap{17}{\PP \VerseOne{1}Unul din cei șapte îngeri care aveau cele șapte cupe a venit și a vorbit cu mine, zicând: “Vino aici. Îți voi arăta judecata marii prostituate care șade pe ape multe,

\VS{2}cu care regii pământului au făcut desfrâu sexual. Cei care locuiesc pe pământ s-au îmbătat cu vinul imoralității ei sexuale.”

\VS{3}El m-a dus în Duhul Sfânt într-un pustiu. Am văzut o femeie șezând pe o fiară de culoare stacojie, plină de nume blasfemiatoare, având șapte capete și zece coarne.

\VS{4}Femeia era îmbrăcată în purpură și stacojiu și împodobită cu aur, pietre prețioase și perle, având în mână un pahar de aur plin de urâciuni și de impuritățile imoralității sexuale de pe pământ.

\VS{5}Și pe fruntea ei era scris un nume: “MISTERUL, BABILONUL CEL MARE, MAMA PROSTITUATEI ȘI A ABOMINAȚIILOR PĂMÂNTULUI.”

\VS{6}Am văzut-o pe femeie îmbătată cu sângele sfinților și cu sângele martirilor lui Isus. Când am văzut-o, m-am mirat cu mare uimire.


\par }{\PP \VS{7}Îngerul mi-a zis: “De ce te miri? Îți voi spune taina femeii și a fiarei care o poartă, care are șapte capete și zece coarne.

\VS{8}Fiara pe care ai văzut-o a fost și nu mai este; și este pe cale să iasă din abis și să meargă la pieire. Cei care locuiesc pe pământ și ale căror nume nu au fost scrise în cartea vieții de la întemeierea lumii se vor minuna când vor vedea că fiara a fost și nu este, și va fi prezentă.
\FTNT{*}{{\FR 17:8 }{\FT{TR citește “încă este” în loc de “va fi prezent”}}}


\par }{\PP \VS{9}Aici este mintea care are înțelepciune. Cele șapte capete sunt șapte munți pe care stă femeia.

\VS{10}Ele sunt șapte împărați. Cinci au căzut, unul este, iar celălalt nu a venit încă. Când va veni, trebuie să mai dureze puțin.

\VS{11}Fiara care era și nu mai este, este și ea însăși a opta, și este dintre cei șapte; și merge la pieire.

\VS{12}Cele zece coarne pe care le-ai văzut sunt zece împărați care nu au primit încă o împărăție, dar primesc autoritate ca împărați împreună cu fiara pentru o oră.

\VS{13}Aceștia au o singură minte și își dau puterea și autoritatea fiarei.

\VS{14}Aceștia se vor lupta împotriva Mielului, iar Mielul îi va birui, pentru că el este Domnul domnilor și Împăratul împăraților; iar cei care sunt cu el sunt chemați, aleși și credincioși.”

\VS{15}El mi-a zis: “Apele pe care le-ai văzut, unde șade prostituata, sunt popoare, mulțimi, națiuni și limbi.

\VS{16}Cele zece coarne pe care le-ai văzut, ele și fiara vor urî prostituata, o vor pustii, o vor dezbrăca, îi vor mânca carnea și o vor arde cu totul în foc.

\VS{17}Căci Dumnezeu le-a pus în inimă să facă ceea ce are în minte, să fie de aceeași părere și să dea împărăția lor fiarei, până când se vor împlini cuvintele lui Dumnezeu.

\VS{18}Femeia pe care ai văzut-o este marea cetate care domnește peste regii pământului.”



\par }\Chap{18}{\PP \VerseOne{1}După acestea, am văzut coborându-se din cer un alt înger, care avea o mare putere. Pământul a fost luminat de slava lui.

\VS{2}El a strigat cu un glas puternic, zicând: “A căzut, a căzut Babilonul cel mare, și a devenit o locuință a demonilor, o închisoare pentru orice duh necurat și o închisoare pentru orice pasăre necurată și urâtă!

\VS{3}Căci toate națiunile au băut din vinul mâniei desfrânării ei, împărații pământului au făcut desfrânare cu ea, iar negustorii pământului s-au îmbogățit din belșugul luxului ei.”


\par }{\PP \VS{4}Și am auzit un alt glas din cer, care zicea: “Ieșiți din ea, poporul Meu, ca să nu luați parte la păcatele ei și să nu vă împărtășiți din plăgile ei,

\VS{5}căci păcatele ei au ajuns până la cer și Dumnezeu și-a adus aminte de nelegiuirile ei.

\VS{6}Întoarceți-vă la ea așa cum s-a întors și răsplătiți-i dublu față de cum a făcut ea și după faptele ei. În paharul pe care l-a amestecat, amestecă-i dublu.

\VS{7}Oricât de mult s-a proslăvit și s-a desfrânat, pe cât de mult dă-i de chin și de jale. Căci ea zice în inima ei: “Șed ca o regină, nu sunt văduvă, și nu voi vedea nicidecum doliu”.

\VS{8}De aceea, într-o singură zi vor veni plăgile ei: moarte, jale și foamete; și va fi arsă cu totul în foc, căci Domnul Dumnezeu, care a judecat-o, este puternic.


\par }{\PP \VS{9}Împărații pământului, care au făcut desfrânare și au trăit în desfrânare cu ea, vor plânge și se vor tângui pe ea, când vor vedea fumul arderii ei,

\VS{10}stând departe de frica chinului ei, și vor zice: “Vai, vai, cetatea cea mare, Babilon, cetatea cea tare! Căci judecata ta a venit într-un ceas!”.

\VS{11}Negustorii pământului plâng și se jelesc din cauza ei, căci nimeni nu le mai cumpără mărfurile:

\VS{12}mărfurile de aur, argint, pietre prețioase, perle, in subțire, purpură, mătase, stacojiu, tot lemnul scump, orice vas de fildeș, orice vas de lemn foarte prețios, de aramă, de fier și de marmură;

\VS{13}și scorțișoară, tămâie, parfum, tămâie, vin, ulei de măsline, făină fină, grâu, oi, cai, care, trupuri și suflete de oameni.

\VS{14}roadele după care a poftit sufletul tău au fost pierdute pentru tine. Toate lucrurile care erau delicate și somptuoase au pierit de la tine și nu le vei mai găsi deloc.

\VS{15}Negustorii acestor lucruri, care s-au îmbogățit prin ea, vor sta departe, de frica chinului ei, plângând și jelind,

\VS{16}spunând: “Vai, vai, cetatea cea mare, cea care era îmbrăcată în in subțire, purpură și stacojiu și împodobită cu aur, pietre prețioase și perle!

\VS{17}Căci într-un ceas, bogății atât de mari sunt pustiite!”. Toți căpitanii de corăbii, toți cei care navighează oriunde, marinarii și toți cei care își câștigă existența pe mare, stăteau departe,

\VS{18}și strigau, când se uitau la fumul arderii ei, zicând: 'Ce este ca cetatea cea mare?

\VS{19}Și-au aruncat țărână pe cap și au strigat, plângând și jelind, zicând: “Vai, vai, cetatea cea mare, în care toți cei care își aveau corăbiile pe mare s-au îmbogățit din cauza marii ei bogății!”. Căci ea a fost pustiită într-un ceas.


\par }{\PP \VS{20}“Bucurați-vă pentru ea, cerule, sfinților, apostolilor și profeților, căci Dumnezeu a judecat-o cu judecata voastră.”


\par }{\PP \VS{21}Un înger puternic a ridicat o piatră ca o piatră mare de moară și a aruncat-o în mare, zicând: “Așa va fi dărâmat Babilonul, cetatea cea mare, și nu va mai fi găsit deloc.

\VS{22}Glasul harpiștilor, al menestreliștilor, al flautiștilor și al trâmbițașilor nu se va mai auzi deloc în voi. Nici un meșter de orice meserie nu va mai fi găsit deloc în voi. Sunetul unei mori nu se va mai auzi deloc în voi.

\VS{23}Lumina unei lămpi nu va mai străluci deloc în voi. Glasul mirelui și al miresei nu se va mai auzi deloc în voi, căci negustorii voștri au fost căpeteniile pământului, căci cu vrăjitoria voastră au fost înșelate toate neamurile.

\VS{24}În ea s-a găsit sângele profeților și al sfinților și al tuturor celor uciși pe pământ.”



\par }\Chap{19}{\PP \VerseOne{1}După acestea, am auzit în ceruri un glas ca un glas tare, ca de mulțime mare, care zicea: “Aleluia! Mântuirea, puterea și gloria aparțin Dumnezeului nostru;

\VS{2}pentru că judecățile lui sunt adevărate și drepte. Căci el a judecat-o pe marea prostituată care a corupt pământul cu imoralitatea ei sexuală și a răzbunat din mâna ei sângele slujitorilor săi.”


\par }{\PP \VS{3}Al doilea a zis: “Aleluia! Fumul ei se ridică în vecii vecilor”.

\VS{4}Cei douăzeci și patru de bătrâni și cele patru făpturi vii au căzut la pământ și s-au închinat lui Dumnezeu, care șade pe tron, spunând: “Amin! Aleluia!”


\par }{\PP \VS{5}Un glas a venit din scaunul de domnie și a zis: “Lăudați pe Dumnezeul nostru, voi toți robii Lui, cei ce vă temeți de El, cei mici și cei mari!”


\par }{\PP \VS{6}Și am auzit ca un glas de mulțime mare, ca un glas de ape multe și ca un glas de tunete puternice, care zicea: “Aleluia! Căci Domnul Dumnezeul nostru, Atotputernicul, domnește!

\VS{7}Să ne bucurăm și să ne veselim nespus de mult și să-i dăm slava. Căci a venit nunta Mielului, iar soția lui s-a pregătit.”

\VS{8}I s-a dat să se îmbrace în in strălucitor, curat și fin, căci inul fin reprezintă faptele drepte ale sfinților.


\par }{\PP \VS{9}El mi-a zis: “Scrie: “Fericiți cei invitați la cina de nuntă a Mielului”.” El mi-a zis: “Acestea sunt cuvinte adevărate ale lui Dumnezeu”.


\par }{\PP \VS{10}M-am aruncat la picioarele Lui, ca să mă închin înaintea Lui. El mi-a zis: “Privește! Nu o face! Eu sunt tovarăș de robie cu tine și cu frații tăi care țin mărturia lui Isus. Adorați-l pe Dumnezeu, căci mărturia lui Isus este Spiritul Profeției”.


\par }{\PP \VS{11}Am văzut cerul deschis și iată un cal alb, și Cel ce ședea pe el se numea Credincios și Adevărat. Cu dreptate el judecă și face război.

\VS{12}Ochii lui sunt o flacără de foc, iar pe capul lui sunt multe coroane. El are nume scrise și un nume scris pe care nimeni nu-l cunoaște decât el însuși.

\VS{13}El este îmbrăcat într-o haină stropită cu sânge. Numele lui se numește “Cuvântul lui Dumnezeu”.

\VS{14}Armatele care sunt în ceruri, îmbrăcate în in alb, pur și fin, îl urmau pe cai albi.

\VS{15}Din gura lui iese o sabie ascuțită, cu două tăișuri, ca să lovească cu ea națiunile. El le va conduce cu un toiag de fier.\FTNT{*}{{\FR 19:15 }{\FT{{\CHAR{Psalmul 2:9}}
}}} El calcă în presa de vin a aprinderii mâniei lui Dumnezeu, Atotputernicul.

\VS{16}El are pe haina și pe coapsa sa un nume scris: “REGELE REGILOR ȘI DOMNUL DOMNILOR”.


\par }{\PP \VS{17}Am văzut un înger care stătea în soare. El a strigat cu glas tare, spunând tuturor păsărilor care zboară pe cer: “Veniți! Adunați-vă la marea cină a lui Dumnezeu,
\FTNT{†}{{\FR 19:17 }{\FT{TR citește “cina marelui Dumnezeu” în loc de “marea cină a lui Dumnezeu”}}}

\VS{18}ca să mâncați carnea regilor, carnea căpeteniilor, carnea vitejilor, carnea cailor și a celor care stau pe ei, carnea tuturor oamenilor, liberi și sclavi, mici și mari!”.

\VS{19}Am văzut fiara, regii pământului și oștile lor, adunate pentru a face război împotriva celui care ședea pe cal și împotriva armatei lui.

\VS{20}Fiara a fost prinsă și, împreună cu ea, falsul profet care făcea semnele pe care le făcea în fața lui, prin care îi înșela pe cei care primiseră semnul fiarei și pe cei care se închinau chipului ei. Aceștia doi au fost aruncați de vii în lacul de foc care arde cu sulf.

\VS{21}Restul au fost uciși cu sabia celui care ședea pe cal, sabia care ieșea din gura lui. Astfel, toate păsările s-au săturat de carnea lor.



\par }\Chap{20}{\PP \VerseOne{1}Am văzut un înger coborând din cer, având în mână cheia prăpastiei și un lanț mare.

\VS{2}El a prins balaurul, șarpele cel vechi, care este diavolul și Satana, care înșală tot pământul locuit\FTNT{*}{{\FR 20:2 }{\FT{TR și NU omit “care amăgește tot pământul locuit”.}}}, l-a legat pentru o mie de ani,

\VS{3}l-a aruncat în abis, l-a închis și l-a pecetluit peste el, ca să nu mai înșele neamurile până la sfârșitul celor o mie de ani. După aceasta, el trebuie să fie eliberat pentru o scurtă perioadă de timp.


\par }{\PP \VS{4}Am văzut niște tronuri, pe care ședeau și pe care se judeca. Am văzut sufletele celor care fuseseră decapitați pentru mărturia lui Isus și pentru cuvântul lui Dumnezeu și ale celor care nu se închinaseră fiarei și chipului ei și nu primiseră semnul de pe frunte și de pe mână. Aceștia au trăit și au domnit cu Hristos timp de o mie de ani.

\VS{5}Ceilalți morți nu au trăit până la sfârșitul celor o mie de ani. Aceasta este prima înviere.

\VS{6}Binecuvântat și sfânt este cel care are parte de prima înviere. Asupra acestora, moartea a doua nu are nicio putere, ci ei vor fi preoți ai lui Dumnezeu și ai lui Hristos și vor domni cu el o mie de ani.


\par }{\PP \VS{7}După cei o mie de ani, Satana va fi eliberat din temnițalui

\VS{8}și va ieși să înșele neamurile care sunt în cele patru colțuri ale pământului, pe Gog și pe Magog, ca să le adune la război, al căror număr este ca nisipul mării.

\VS{9}Ei s-au suit pe lățimea pământului și au înconjurat tabăra sfinților și cetatea iubită. Focul s-a coborât din cer de la Dumnezeu și i-a mistuit.

\VS{10}Diavolul care i-a înșelat a fost aruncat în iazul de foc și de sulf, unde sunt și fiara și falsul profet. Ei vor fi chinuiți zi și noapte în vecii vecilor.


\par }{\PP \VS{11}Am văzut un tron mare și alb și pe Cel ce ședea pe el, de la fața căruia fugeau pământul și cerul. Nu s-a găsit loc pentru ei.

\VS{12}Am văzut pe cei morți, mari și mici, stând în picioare înaintea tronului, și au deschis cărți. S-a deschis o altă carte, care este cartea vieții. Morții au fost judecați din cele scrise în cărți, după faptele lor.

\VS{13}Marea a dat pe față morții care se aflau în ea. Moartea și Hadesul\FTNT{†}{{\FR 20:13 }{\FT{sau, Iadul
}}} au renunțat la morții care erau în ele. Au fost judecați, fiecare după faptele sale.

\VS{14}Moartea și Hadesul\FTNT{‡}{{\FR 20:14 }{\FT{sau, Iadul}}} au fost aruncate în lacul de foc. Aceasta este a doua moarte, lacul de foc.

\VS{15}Dacă cineva nu era găsit scris în cartea vieții, era aruncat în lacul de foc.



\par }\Chap{21}{\PP \VerseOne{1}Am văzut un cer nou și un pământ nou, căci cerul dintâi și pământul dintâi au trecut și marea nu mai este.

\VS{2}Am văzut cetatea sfântă, Noul Ierusalim, coborând din cer de la Dumnezeu, pregătită ca o mireasă împodobită pentru soțul ei.

\VS{3}Am auzit un glas puternic din cer care zicea: “Iată, locuința lui Dumnezeu este cu oamenii; El va locui cu ei, iar ei vor fi poporul Lui și Dumnezeu însuși va fi cu ei ca Dumnezeul lor.

\VS{4}El va șterge orice lacrimă din ochii lor. Moartea nu va mai fi; nu va mai fi nici doliu, nici strigăt, nici durere. Cele dintâi lucruri au trecut.”


\par }{\PP \VS{5}Cel ce șade pe scaunul de domnie a zis: “{\WJ{Iată, Eu fac toate lucrurile noi.”}} El a zis:
{\WJ{“Scrieți, căci aceste cuvinte ale lui Dumnezeu sunt credincioase și adevărate.”}}

\VS{6}El mi-a spus:
{\WJ{“Eu sunt Alfa și Omega, Începutul și Sfârșitul. Eu voi da gratuit celui care este însetat de la izvorul apei vieții.
}}

\VS{7}Celui care va
{\WJ{birui, îi voi da aceste lucruri. Eu voi fi Dumnezeul lui, iar el va fi fiul meu.
}}

\VS{8}{\WJ{Dar pentru cei lași, necredincioși, păcătoși,
}}\FTNT{*}{{\FR 21:8 }{\FT{TR și NU omit “păcătoși”}}}{\WJ{abominabili, ucigași, imorali, vrăjitori, închinători
}}\FTNT{†}{{\FR 21:8 }{\FT{Cuvântul pentru “vrăjitori” include aici și pe cei care folosesc poțiuni și droguri.}}}{\WJ{la idoli și toți mincinoșii, partea lor este în lacul care arde cu foc și sulf, care este moartea a doua.”
}}


\par }{\PP \VS{9}Unul din cei șapte îngeri care aveau cele șapte cupe încărcate cu cele șapte plăgi de pe urmă a venit și a vorbit cu mine, zicând: “Vino aici. Îți voi arăta mireasa, soția Mielului”.

\VS{10}El m-a dus în Duhul Sfânt pe un munte mare și înalt și mi-a arătat cetatea sfântă, Ierusalimul, care se cobora din cer de la Dumnezeu,

\VS{11}având gloria lui Dumnezeu. Lumina ei era ca o piatră foarte prețioasă, ca o piatră de iaspis, limpede ca cristalul;

\VS{12}având un zid mare și înalt, cu douăsprezece porți și, la porți, doisprezece îngeri și nume scrise pe ele, care sunt numele celor douăsprezece triburi ale copiilor lui Israel.

\VS{13}La răsărit erau trei porți, la nord trei porți, la sud trei porți, și la vest trei porți.

\VS{14}Zidul cetății avea douăsprezece temelii și pe ele douăsprezece nume ale celor doisprezece apostoli ai Mielului.


\par }{\PP \VS{15}Cel care vorbea cu mine avea ca măsură o trestie de aur, ca să măsoare cetatea, porțile și zidurile ei.

\VS{16}Cetatea este pătrată. Lungimea lui este la fel de mare ca lățimea lui. El a măsurat cetatea cu trestia: douăsprezece mii doisprezece stadii.\FTNT{‡}{{\FR 21:16 }{\FT{12.012 stadii = sau 2.221 kilometri sau 1.380 mile. TR citește 12.000 de stadii în loc de 12.012 stadii.}}} Lungimea, lățimea și înălțimea ei sunt egale.

\VS{17}Zidul ei este de o sută patruzeci și patru de coți,
\FTNT{§}{{\FR 21:17 }{\FT{144 coți înseamnă aproximativ 65,8 metri sau 216 picioare.}}}după măsura unui om, adică a unui înger.

\VS{18}Construcția zidului ei era de jaspis. Orașul era de aur pur, ca sticla pură.

\VS{19}Temeliile zidului cetății erau împodobite cu tot felul de pietre prețioase. Prima temelie era de iaspis, a doua de safir,\FTNT{*}{{\FR 21:19 }{\FT{sau lapis lazuli}}} a treia de calcedonie, a patra de smarald,

\VS{20}a cincea de sardonyx, a șasea de sardiu, a șaptea de crisolit, a opta de beril, a noua de topaz, a zecea de crioprasă, a unsprezecea de jacint și a douăsprezecea de ametist.

\VS{21}Cele douăsprezece porți erau douăsprezece perle. Fiecare dintre porți era făcută dintr-o perlă. Strada orașului era de aur pur, ca o sticlă transparentă.


\par }{\PP \VS{22}Nu am văzut în ea nici un templu, căci Domnul Dumnezeu Atotputernic și Mielul sunt templul ei.

\VS{23}Cetatea nu are nevoie de soare sau de lună pentru a străluci, căci însăși slava lui Dumnezeu a luminat-o, iar lampa ei este Mielul.

\VS{24}Națiunile vor umbla în lumina ei. Regii pământului aduc în ea gloria și onoarea națiunilor.

\VS{25}Porțile ei nu vor fi în niciun fel închise ziua (pentru că acolo nu va fi noapte),

\VS{26}și ei vor aduce gloria și onoarea națiunilor în ea, ca să poată intra.

\VS{27}Nu va intra în ea în niciun fel ceva profanator sau cineva care provoacă o urâciune sau o minciună, ci numai cei care sunt scriși în cartea vieții Mielului.



\par }\Chap{22}{\PP \VerseOne{1}Mi-a arătat un
\FTNT{*}{{\FR 22:1 }{\FT{TR adaugă “pur”}}}râu de apă a vieții, limpede ca cristalul, care ieșea din scaunul de domnie al lui Dumnezeu și al Mielului,

\VS{2}în mijlocul străzii lui. De o parte și de alta a râului era pomul vieții, care purta douăsprezece feluri de fructe și care își dădea rodul în fiecare lună. Frunzele copacului erau pentru vindecarea națiunilor.

\VS{3}Nu va mai fi niciun blestem. Tronul lui Dumnezeu și al Mielului va fi în ea, iar slujitorii lui îi vor sluji.

\VS{4}Îi vor vedea fața și numele lui va fi pe frunțile lor.

\VS{5}Nu va mai fi noapte și nu vor avea nevoie de lumina lămpii sau de lumina soarelui, pentru că Domnul Dumnezeu îi va lumina. Ei vor domni în vecii vecilor.


\par }{\PP \VS{6}El mi-a zis: “Cuvintele acestea sunt credincioase și adevărate. Domnul Dumnezeul spiritelor profeților a trimis pe îngerul Său ca să arate robilor Săi lucrurile care trebuie să se întâmple curând.”


\par }{\PP \VS{7}“{\WJ{Iată, vin curând! Binecuvântat este cel care păzește cuvintele profeției din cartea aceasta.”
}}


\par }{\PP \VS{8}Și eu, Ioan, sunt cel ce a auzit și a văzut aceste lucruri. Când am auzit și am văzut, am căzut să mă închin la picioarele îngerului care mi-a arătat aceste lucruri.

\VS{9}El mi-a spus: “Nu trebuie să faci asta! Eu sunt tovarăș de robie cu tine și cu frații tăi, cu profeții și cu cei care păzesc cuvintele acestei cărți. Adoră-l pe Dumnezeu”.

\VS{10}El mi-a zis: “Nu pecetlui cuvintele profeției acestei cărți, căci timpul este aproape.

\VS{11}Cine acționează nedrept, să acționeze tot nedrept. Cine se murdărește, să se murdărească în continuare. Cine este neprihănit, să facă încă dreptate. Cine este sfânt, să fie încă sfânt.”


\par }{\PP \VS{12}{\WJ{“Iată, vin curând! Răsplata Mea este cu Mine, ca să răsplătesc fiecăruia după fapta lui.
}}

\VS{13}{\WJ{Eu sunt Alfa și Omega, Cel dintâi și Cel de pe urmă, Începutul și Sfârșitul.
}}

\VS{14}{\WJ{Ferice de cei care împlinesc poruncile Lui,
}}\FTNT{†}{{\FR 22:14 }{\FT{NU citește “își spală hainele” în loc de “împlinesc poruncile Lui”.}}}{\WJ{ca să aibă dreptul la pomul vieții și să intre pe porțile cetății.
}}

\VS{15}{\WJ{Afară sunt câinii, vrăjitorii, cei imorali, ucigașii, închinătorii la idoli și toți cei care iubesc și practică minciuna.
}}

\VS{16}{\WJ{Eu, Isus, am trimis pe îngerul meu să vă mărturisească aceste lucruri pentru adunări. Eu sunt rădăcina și urmașul lui David, Strălucirea și Steaua dimineții.”
}}


\par }{\PP \VS{17}Duhul și mireasa spun: “Vino!” Cel care aude, să spună: “Vino!” Cine este însetat, să vină! Cine dorește, să ia gratuit apa vieții.


\par }{\PP \VS{18}Eu mărturisesc oricui ascultă cuvintele profeției din cartea aceasta că, dacă cineva adaugă la ele, Dumnezeu îi va adăuga plăgile care sunt scrise în cartea aceasta.

\VS{19}Dacă cineva va tăia din cuvintele cărții acestei profeții, Dumnezeu îi va lua partea lui din pomul\FTNT{‡}{{\FR 22:19 }{\FT{TR citește “Carte” în loc de “copac”}}} vieții și din cetatea sfântă, care sunt scrise în cartea aceasta.

\VS{20}Cel care mărturisește aceste lucruri spune:
{\WJ{“Da, voi veni curând”.
}}

\par }{\PP Amin! Da, vino, Doamne Isuse!


\par }{\PP \VS{21}Harul Domnului Isus Hristos să fie cu toți sfinții. Amin.


\par }